\section{LaMET 原理：内容、几种表述与地位}\label{sec:lamet-principle}

出发点是与重夸克有效理论的类比：取 $m_Q\to\infty$ 只改变紫外物理，差别是可微扰计算的匹配系数。季向东提出，对以动量 $P^z$ 运动的强子，“先算等时关联函数、再令 $P^z\to\infty$”与“先令 $P^z\to\infty$（光前量子化）、再算关联函数”也只差紫外物理\cite{Ji2013,Ji:2014gla}。形式上的无穷推动极限相同的一族等时算符称为一个\emph{普适类}\cite{Ji:2014gla,Hatta:2013gta}：$\Gamma=\gamma^0$ 与 $\gamma^z$ 的准 PDF 算符都趋向 $\gamma^+$ 的光线算符（\cref{def:pdf-lightray}），匹配核不同（\cref{prop:lamet-oneloop}），极限相同。

\begin{conjecture}[LaMET 原理（非微扰形式） \tagC]\label{conj:lamet-principle}
设 $\mathcal O_E$ 为一个重整化的等时（欧氏）算符，其形式上的无穷推动极限是光前算符 $\mathcal O_{LF}$（记 $\mathcal O_E\in[\mathcal O_{LF}]$；例如 $\Gamma=\gamma^0$ 与 $\gamma^z$ 的准 PDF 算符都属于 $\gamma^+$ 的普适类，其光前算符给出\cref{def:pdf-quark}的夸克分布）。则对 $P^z\gg\Lqcd,M$ 的强子态，
\begin{equation}\label{eq:lamet-principle}
\begin{split}
\langle P|\mathcal O_E|P\rangle(P^z,\mu)={}&\Big[\mathcal Z\Big(\frac{P^z}{\mu},\as(\mu)\Big)\otimes\langle P|\mathcal O_{LF}|P\rangle(\mu)\Big]\\
&+\sum_{k\ge1}\Big(\frac{\Lambda}{P^z}\Big)^{2k}\times\big(\text{高扭度等时算符的矩阵元}\big),
\end{split}
\end{equation}
其中：$\mathcal Z$ 有不含红外敏感性的微扰展开；幂次展开是渐近展开，按部分子的硬标度（$|x|P^z$、$|1-x|P^z$）而不是强子动量组织；当 Wilson 线自能与匹配系数用同一个重整子处方定义时，不出现 $\Lambda/P^z$ 的线性项（见本章后文对线性重整子的讨论）。

非微扰地位：\cref{asm:intro-wightman}、\cref{asm:intro-composite}、\cref{asm:intro-af-coeff}（MA0、MA4）。
\uses{def:pdf-quark, def:pdf-lightray}
\end{conjecture}

\paragraph{地位与缺口}
对任何一个普适类的任何一个成员，都没有非微扰的证明；证明需要 QCD 中的非微扰 OPE 或非微扰因子化定理，而这两者本身都是猜想（\cref{conj:ee-np-ope}、\cref{asm:intro-np-ope}；参见\cref{rem:power-np-status}）。\cref{eq:lamet-principle}的微扰内容——部分子态、领头幂次、逐阶——是本章末的\cref{cor:lamet-principle-pt}（\tagA）；共线情形的非微扰内容被精确地归结为三个编号假设\cref{asm:lamet-uniform,asm:lamet-twist4,asm:lamet-large-z}（\tagC），即\cref{thm:lamet-hadronic}的前提。无线性项的说法是\cref{prop:lamet-lrr}（\tagP，大 $\beta_0$ 陈述）的内容。支持这一猜想的证据（不是证明）有：重整子分析（\cref{prop:lamet-bvz}）、二维可严格求解模型（\cref{thm:yzliu-gn-convergence}，第~23~章）与格点数据。季向东所说“由于 QCD 的渐近自由，两种极限次序只在微扰论中不同，对非微扰物理等价”\cite{Ji:2024oka}是一个启发式论证：它没有说明非微扰的余项如何控制。文献：\cite{Ji2013,Ji:2014gla}，\cite[第~II--III~节]{JiEtAl2021}，\cite{Ji:2020byp,Ji:2024oka}。要把标签提高到 \tagA，至少需要一个对强子态成立、余项关于 Ioffe 时间一致的非微扰算符乘积展开（\cref{asm:lamet-uniform}的解除途径）。

\begin{remark}[LaMET 原理的几种表述及其可证明的内容]\label{rem:lamet-formulations}
(a)~\emph{矩的形式}\cite{Ji2013}。扭度 2 算符 $O^{\mu_1\cdots\mu_n}=\mathcal S\,\bar\psi\gamma^{\mu_1}\ii D^{\mu_2}\cdots\ii D^{\mu_n}\psi-\text{迹}$ 满足 $\langle P|O^{\mu_1\cdots\mu_n}|P\rangle=2a_n(P^{\mu_1}\cdots P^{\mu_n}-\text{迹})$，$a_n=\int_{-1}^1\dd x\,x^{n-1}f(x)$（\cref{thm:pdf-moments-local}）。全部指标取 $z$ 分量时迹项可精确求出（\cref{prop:lamet-tmc}的证明）：
\begin{equation}\label{eq:lamet-ji-moments}
\begin{gathered}
\langle P|\bar\psi\gamma^z(\ii D^z)^{n-1}\psi|P\rangle=2a_n(P^z)^nK_n(c)+(\text{高扭度}),\\
K_n(c)=\sum_{i\ge0}\binom{n-i}{i}c^i=1+(n-1)c+\order{c^2},\qquad c=\frac{M^2}{4P_z^2}.
\end{gathered}
\end{equation}
这是 Lorentz 协变性的推论（严格；\cite{Ji2013}写作 $2a_n(P^z)^n[1+\order{\Lambda^2/P_z^2,M^2/P_z^2}]$）。由此“反演”到非定域算符却是启发式的：对 $n$ 求和与 $P^z\to\infty$ 不可交换（$\MSbar$ 准 PDF 没有矩，\cref{prop:lamet-tails}）；紫外重整化与 $P^z\to\infty$ 也不可交换——\cite{Ji2013}在 $A^z=0$ 规范中用横动量截断 $\mu$ 算出的单圈准 PDF 含 $\Lambda_\mu/P^z$，$\Lambda_\mu=\sqrt{\mu^2+(1-x)^2P_z^2}-(1-x)P^z$：先取 $\mu\to\infty$ 得线性发散，先取 $P^z\to\infty$ 才得 DGLAP 对数。
(b)~\emph{有效场论形式}\cite{Ji:2014gla}：$F(P^z/\Lambda_{\mathrm{UV}})=Z(P^z/\Lambda_{\mathrm{UV}},\Lambda_{\mathrm{UV}}/\mu)\,f(\mu)+\order{\Lambda^2/P_z^2}$，$\gamma=\partial\ln Z/\partial\ln P^z$；该文引入普适类（$\Delta G$ 见\cite{Hatta:2013gta}，据其摘要）、准光前波函数与最小可分辨动量分数 $x_{\min}\sim\Lambda/P^z$，并断言（未给证明）“Feynman 图层次的因子化可以逐阶证明”。
(c)~\emph{动量密度与动量重整化群}\cite[第~II--III~节]{JiEtAl2021}：准 PDF 被诠释为夸克的动量密度 $N_P(k^z)$；由\cref{prop:lamet-oneloop}的 $\MSbar$ 核与 $\dd\ell_p/\dd\ln P^z=-2$（\cref{prop:lamet-oneloop}(f)）得单圈方程
\begin{equation}\label{eq:lamet-momentum-rge}
P^z\frac{\dd\tilde q}{\dd P^z}=\frac{\as C_F}{\pi}\Big[\Big[\frac{1+y^2}{1-y}\Big]_+-\frac32\delta(1-y)\Big]\otimes f+\order{\as^2};
\end{equation}
该文断言匹配公式及其逆对所有阶成立。
(d)~\emph{“Weinberg 意义下”的欧氏有效理论}\cite{Ji:2020byp,Ji:2024oka}：共线夸克的领头大速度拉氏量 $\mathcal L^{(0)}_{q,v}=\bar\psi_v\big[\ii v\cdot D+\frac{\ii\bar v\cdot D}{2\gamma^2}+\ii D_\perp\frac1{2\ii\bar v\cdot D}\ii D_\perp\big]\sla{\bar v}\psi_v$（$v^\mu=(1,0,0,v)$，$v^2=1/\gamma^2$，$\bar v^\mu=(v,0,0,-1)/(2v)$，$v\cdot\bar v=1$；拉氏量只通过协变导数含耦合常数，在本书约定 $D=\partial-\ii gA$ 下形式不变）；有限推动是快度发散的非微扰正规子；扭度 4 等时关联函数的幂次发散须与领头系数的红外重整子一致地定义（类比 Symanzik 改进）；展开在 $x\sim\Lambda/P^z$ 处失效；有限 $P^z$ 下准 PDF 在 $y=0,1$ 处解析（精确形式见\cref{lem:lamet-strip}）。
四种表述中只有 (a) 的矩关系与 (c) 中单圈方程的推导是严格的；“所有阶成立”由\cref{sec:lattice-matching}的匹配定理（\tagA）给出全阶论证，非微扰的断言即\cref{conj:lamet-principle}。
\end{remark}
